\documentclass[12pt]{article}
\usepackage[utf8]{inputenc}
\usepackage[T1]{fontenc}
\usepackage{hyperref}
\title{Rom Language Resources for K-12 Curriculum Development}
\author{Compiled by Claude Sonnet 5 (claude-sonnet-5) for Andr\'as Kornai}
\date{}
\begin{document}
\maketitle
\section{Introduction}
Rom (ISO 639-3 tdk; Glottocode tamb1267; also called Tambas) is described as an Afro-Asiatic (West Chadic, Ron group) language of Plateau State, Nigeria, with about 3,000 speakers. It must not be confused with Horom (ISO 639-3 hoe), which is a different, Plateau (Niger--Congo) language. Search results were ambiguous about the name Rom, and the sources do not settle how it relates to Tambas or to Horom. Very little was found.
\section{Primary Educational Resources}
None found.
\section{Linguistic and Reference Materials}
\begin{itemize}
\item Ron languages (Wikipedia; lists Tambas in the Ron branch): \url{https://en.wikipedia.org/wiki/Ron_languages}
\item Numerals page for ``Rom, Plateau state, Nigeria'' (also named Marom, Tambas; the page gives about 7,000 speakers): \url{https://lingweb.eva.mpg.de/channumerals/Rom.htm}
\end{itemize}
\section{Religious and Community Resources}
None found.
\section{Digital and Audio Resources}
A Ron-languages Webonary page appeared in results but was not examined, so its relevance to Rom is unverified: \url{https://www.webonary.work/ron/?p=9444}
\section{Curriculum Development Context}
Nothing substantial was found. Any work would first need clarification of the language's identity and documentation of it.
\end{document}